% Part: methods
% Chapter: proofs
% Section: cant-do-it

\documentclass[../../../include/open-logic-section]{subfiles}

\begin{document}

\olfileid[hi]{mth}{prf}{cnt}

\olsection{मैं यह नहीं कर सकता!{}}

हम सभी ऐसे बिंदु पर पहुँचते हैं जहाँ हार मानने का मन होता है। किंतु आप यह
कर \emph{सकते हैं}। आपके शिक्षक, शिक्षण सहायक और सहपाठी सहायता कर सकते हैं।
उनसे सहायता माँगिए!{} संकट से बचने और हार मानने का मन होने पर क्या करें,
इसके कुछ सुझाव ये हैं।

यह सुनिश्चित करने के लिए कि आप समस्याएँ सफलतापूर्वक हल कर सकें, निम्न करें:
\begin{enumerate}
\item \emph{यथासंभव पहले आरंभ करें।} पूरे सत्र में हम व्यस्त होते हैं और
  हममें से अनेक टालमटोल से जूझते हैं; सबसे अच्छे कामों में से एक है गृहकार्य
  जल्दी आरंभ करना। इस प्रकार यदि अटक जाएँ, तो आपके पास समाधान ढूँढने का समय
  होगा (रोने के अतिरिक्त)।
\item \emph{अपने सहपाठियों से बात करें।} आप अकेले नहीं हैं। कक्षा के अन्य
  लोग भी कठिनाई में हो सकते हैं---पर उनकी कठिनाइयाँ भिन्न हो सकती हैं।
  साथियों से चर्चा समस्या पर ऐसा भिन्न दृष्टिकोण दे सकती है जिससे समाधान
  सूझे। निस्संदेह उनका समाधान केवल नकल न करें: संकेत माँगें, या बताएँ कि कहाँ
  अटके और अगला कदम पूछें। और जब समझ जाएँ, तो बदले में सहायता करें। किसी और
  की सहायता और बातों की व्याख्या आपकी समझ भी बेहतर करेगी।
\item \emph{सहायता माँगें।} आपके लिए अनेक संसाधन उपलब्ध हैं---आपके शिक्षक
  और शिक्षण सहायक आपकी सहायता के लिए हैं और \emph{चाहते हैं} कि आप सफल हों।
  वे समस्या हल करने और प्रक्रिया में आपकी कठिनाई का स्थान पहचानने में
  सहायता कर सकेंगे।
\item \emph{विराम लें।} यदि अटके हैं, तो इसका कारण यह \emph{हो सकता है} कि
  आप समस्या को बहुत देर से देखते रहे हैं। छोटा विराम लें, एक कप चाय पिएँ,
  या कुछ समय दूसरी समस्या पर काम करें, फिर ताज़े मन से लौटें। रात भर सोचने
  दें।
\end{enumerate}

ध्यान दें कि इन रणनीतियों के लिए प्रमाण पर बहुत पहले काम आरंभ करना आवश्यक
है। यदि नियत दिन से पहले रात दो बजे प्रमाण आरंभ किया, तो वे अधिक सहायक नहीं
होंगी।

यह बहुत निराशाजनक लग सकता है, किंतु प्रमाण हल करना ऐसी चुनौती है जिसका अंत
में फल मिलता है। कुछ लोग इसे पेशे के रूप में करते हैं---अतः इसमें आनंद की
कोई बात अवश्य होगी। लगभग हर चीज़ की तरह समस्याएँ हल करने और प्रमाण करने के
लिए अभ्यास चाहिए। आपको ऐसे सहपाठी दिख सकते हैं जिन्हें यह सरल लगता है:
संभवतः वे पहले ही बहुत अभ्यास कर चुके हैं। बहुत जल्दी हार न मानें।

यदि किसी विशेष समस्या पर समय (या धैर्य) समाप्त हो जाए, तो कोई बात नहीं।
इसका अर्थ यह नहीं कि आप मूर्ख हैं या कभी समझ नहीं पाएँगे। अपने शिक्षक या
दूसरे विद्यार्थी से जानें कि वह कैसे किया जाता है और पहचानें कि कहाँ गलती
हुई या अटके, ताकि अगली बार समान प्रश्न में उससे बच सकें। फिर समाधान देखे
बिना उसे करने का प्रयास करें। और अगली बार पहले आरंभ करें (और सहायता भी पहले
माँगें)।

\end{document}
